\subsubsection{SWAP gate}

The \definition{SWAP gate} is a \textbf{two-qubit gate} that \hl{exchanges the states of two qubits} and has \hl{no control/target distinction}, it is \textbf{symmetric}. It simply swaps:
\begin{equation}
    \ket{v_{A}} \otimes \ket{v_{B}} \xrightarrow{SWAP} \ket{v_{B}} \otimes \ket{v_{A}}
\end{equation}

\highspace
\begin{flushleft}
    \textcolor{Green3}{\faIcon{book} \textbf{Action on Basis States}}
\end{flushleft}
The SWAP gate acts on the basis states as follows:

\begin{table}[!htp]
    \centering
    \begin{tabular}{@{} c | c @{}}
        \toprule
        \textbf{Input} & \textbf{Output} \\
        \midrule
        $\ket{00}$ & $\ket{00}$ \\
        $\ket{01}$ & $\ket{10}$ \\
        $\ket{10}$ & $\ket{01}$ \\
        $\ket{11}$ & $\ket{11}$ \\
        \bottomrule
    \end{tabular}
\end{table}

\noindent
It essentially \textbf{swaps the second and third basis states}, while leaving the first and fourth unchanged.

\highspace
\begin{flushleft}
    \textcolor{Green3}{\faIcon{square-root-alt} \textbf{Matrix Representation}}
\end{flushleft}
The matrix representation of the SWAP gate in the computational basis $\{\ket{00},\break \ket{01}, \ket{10}, \ket{11}\}$ is:
\begin{equation}
    \text{SWAP} = \begin{bmatrix}
        1 & 0 & 0 & 0 \\
        0 & 0 & 1 & 0 \\
        0 & 1 & 0 & 0 \\
        0 & 0 & 0 & 1
    \end{bmatrix}
\end{equation}
This matrix clearly shows the \textbf{swapping of the second and third basis states}, while the first and fourth remain unchanged.

\highspace
\begin{flushleft}
    \textcolor{Green3}{\faIcon{book} \textbf{Action on General States}}
\end{flushleft}
Let's consider two arbitrary single-qubit states:
\begin{equation*}
    \ket{v_A} = a_0 \ket{0} + a_1 \ket{1}, \quad \ket{v_B} = b_0 \ket{0} + b_1 \ket{1}
\end{equation*}
The tensor product of these states is:
\begin{equation*}
    \ket{v_A} \otimes \ket{v_B} = a_0 b_0 \ket{00} + a_0 b_1 \ket{01} + a_1 b_0 \ket{10} + a_1 b_1 \ket{11}
\end{equation*}
Applying the SWAP gate to this state gives:
\begin{align*}
    \text{SWAP}(\ket{v_A} \otimes \ket{v_B}) &= a_0 b_0 \ket{00} + \text{\hl{$a_1 b_0 \ket{01}$}} + \text{\hl{$a_0 b_1 \ket{10}$}} + a_1 b_1 \ket{11} \\
    &= (b_0 \ket{0} + b_1 \ket{1}) \otimes (a_0 \ket{0} + a_1 \ket{1}) \\
    &= \ket{v_B} \otimes \ket{v_A}
\end{align*}
This confirms that the SWAP gate indeed \textbf{exchanges the states of the two qubits}, as expected.

\highspace
\begin{flushleft}
    \textcolor{Green3}{\faIcon{tools} \textbf{Circuit Representation}}
\end{flushleft}
In quantum circuit diagrams, the SWAP gate is typically represented by two crossed lines with $\times$ symbols at the crossing point:

\begin{center}
    \begin{quantikz}
        \lstick{$\ket{v_{A}}$} & \swap{1} & \rstick{$\ket{v_{B}}$} \\
        \lstick{$\ket{v_{B}}$} & \targX{} & \rstick{$\ket{v_{A}}$}
    \end{quantikz}
\end{center}

\noindent
This notation visually emphasizes the \textbf{swapping of the two qubits' states}.

\highspace
\begin{flushleft}
    \textcolor{Green3}{\faIcon{cogs} \textbf{Properties}}
\end{flushleft}
\begin{itemize}
    \item \important{Unitary}
    \begin{equation*}
        \text{SWAP}^\dagger \text{SWAP} = I
    \end{equation*}
    This means that the SWAP gate preserves the inner product and is reversible.

    \item \important{Reversible (self-inverse)}
    \begin{equation*}
        \text{SWAP}^2 = I
    \end{equation*}
    This means that applying the SWAP gate twice will return the system to its original state.

    \item \important{Symmetric (no control/target)}. The SWAP gate does not distinguish between the two qubits; it treats them symmetrically. This is in contrast to gates like CNOT, which have a clear control and target qubit.
    
    \item \important{Does NOT create entanglement}. The SWAP gate simply exchanges the states of two qubits and does \textbf{not introduce any new correlations} between them. If the input state is separable (not entangled), the output state will also be separable (and vice versa).
    
    \item \important{Can be decomposed into 3 CNOT gates}. The \textbf{SWAP gate can be implemented using three CNOT gates}, which is important for practical implementations on quantum hardware that may not have a native SWAP gate.

    \begin{proof}
        We want to transform:
        \begin{equation*}
            \ket{a, \, b} \rightarrow \ket{b, \, a}
        \end{equation*}
        But CNOT only does:
        \begin{equation*}
            \ket{a, \, b} \rightarrow \ket{a, \, a \oplus b}
        \end{equation*}
        So we need to \textbf{combine multiple CNOTs} to achieve a SWAP. Let's track what happens to a generic input:
        \begin{equation*}
            \ket{a, \, b} \quad \text{with } a, b \in \{0, 1\}
        \end{equation*}
        \begin{itemize}
            \item \textbf{First} CNOT (control: first qubit, target: second qubit):
            \begin{equation*}
                \ket{a, \, b} \xrightarrow{\text{CNOT}} \ket{a, \, a \oplus b}
            \end{equation*}
            \begin{center}
                \begin{quantikz}
                    \lstick{$\ket{a}$} & \ctrl{1} & \rstick{$\ket{a}$} \\
                    \lstick{$\ket{b}$} & \targ{}  & \rstick{$\ket{a \oplus b}$}
                \end{quantikz}
            \end{center}

            \item \textbf{Second} CNOT (control: second qubit, target: first qubit):
            \begin{equation*}
                \ket{a, \, a \oplus b} \xrightarrow{\text{CNOT}} \ket{a \oplus (a \oplus b), \, a \oplus b} = \ket{b, \, a \oplus b}
            \end{equation*}
            \begin{center}
                \begin{quantikz}
                    \lstick{$\ket{a}$} & \ctrl{1} & \push{\ket{a}}            & \targ{}   & \rstick{$\ket{b}$} \\
                    \lstick{$\ket{b}$} & \targ{}  & \push{\ket{a \oplus b}}   & \ctrl{-1} & \rstick{$\ket{a \oplus b}$}
                \end{quantikz}
            \end{center}

            \item \textbf{Third} CNOT (control: first qubit, target: second qubit):
            \begin{equation*}
                \ket{b, \, a \oplus b} \xrightarrow{\text{CNOT}} \ket{b, \, b \oplus (a \oplus b)} = \ket{b, a}
            \end{equation*}
            \begin{center}
                \begin{quantikz}
                    \lstick{$\ket{a}$} & \ctrl{1} & \push{\ket{a}}            & \targ{}   & \push{\ket{b}}          & \ctrl{1}  & \rstick{$\ket{b}$} \\
                    \lstick{$\ket{b}$} & \targ{}  & \push{\ket{a \oplus b}}   & \ctrl{-1} & \push{\ket{a \oplus b}} & \targ{}   & \rstick{$\ket{a}$}
                \end{quantikz}
            \end{center}
        \end{itemize}
        In summary, the sequence of operations is:
        \begin{center}
            \begin{quantikz}
                    \lstick{$\ket{a}$} & \ctrl{1}\gategroup[2,steps=3,style={dashed,rounded corners,fill=blue!20,inner xsep=2pt},background,label style={label position=below,anchor=north,yshift=-0.2cm}]{{\sc swap}} & \targ{} & \ctrl{1} & \rstick{$\ket{b}$} \\
                    \lstick{$\ket{b}$} & \targ{} & \ctrl{-1} & \targ{} & \rstick{$\ket{a}$}
            \end{quantikz}
        \end{center}
    \end{proof}
\end{itemize}
In general, the SWAP gate \textbf{does not create new correlations}; it only \hl{reorderes the information between qubits}. In other words, it only \hl{changes where the information is stored}.